\newcommand{\mvteaserwidthb}{.333\textwidth}
\begin{figure*}[t]
\captionsetup[subfigure]{labelformat=empty}
\centering
\makebox[\textwidth][c]{%
\renewcommand{\arraystretch}{0.4}%
\begin{tabular}{@{}c@{\betweencols}c@{\betweencols}c@{}} % 3 columns
\begin{subfigure}[t]{\mvteaserwidthb}
  \includegraphics[width=\columnwidth]{figures/result_grid_mv/chair.pdf}
\end{subfigure} &
\begin{subfigure}[t]{\mvteaserwidthb}
  \includegraphics[width=\columnwidth]{figures/result_grid_mv/orangeguy.pdf}
\end{subfigure} &
\begin{subfigure}[t]{\mvteaserwidthb}
  \includegraphics[width=\columnwidth]{figures/result_grid_mv/unitree.pdf}
\end{subfigure} \\ 
\end{tabular}%
}
\caption{
\textbf{Multi-view conditioning.} Our decoupled joint modeling naturally supports multi-view conditioning. When multiple input views are available, we fuse the shared \textit{view-agnostic} object latent across flow paths (similar to MultiDiffusion~\citep{bar2023multidiffusion}), enabling object and cameras refinement across all views. Top: inputs; Middle: reconstructed 3D object and camera poses; Bottom: rendered images and geometry.}
\label{fig:mv_images}
\end{figure*}
